Celem niniejszej pracy było zaprojektowanie i~zaimplementowanie interaktywnego
silnika symulacji cieczy, wykonującego pełny potok obliczeniowy --- od dynamiki
cząstek po gotowy obraz --- w~czasie rzeczywistym na procesorze graficznym.
Punkt wyjścia stanowiła analiza eulerowskiego i~lagrange'owskiego opisu ruchu
płynów oraz matematyczne sformułowanie metody SPH wraz z~wariantem DFSPH,
wymuszającym nieściśliwość cieczy przez dwuetapową korekcję gęstości
i~dywergencji. Tak zdefiniowane podstawy teoretyczne przełożono na implementację
w~języku Rust z~wykorzystaniem API Vulkan, a~każdą istotną decyzję projektową
uzasadniono analizą strukturalną lub bezpośrednim pomiarem.

W~ramach pracy zrealizowano następujące zadania:
\begin{itemize}
    \item dokonano przeglądu istniejących rozwiązań symulacji cieczy
          (Blender, Unreal Engine~5, SPlisHSPlasH) i~wyznaczono na ich tle
          miejsce własnego systemu: czysto cząsteczkowej implementacji DFSPH
          wykonywanej w~całości na GPU;
    \item zaimplementowano kompletny solver DFSPH jako zestaw jedenastu potoków
          obliczeniowych, obejmujących wyznaczanie gęstości i~współczynnika
          $\alpha$, iteracyjne pętle korekcji ciśnienia oraz całkowanie ruchu
          cząstek;
    \item opracowano wyszukiwanie sąsiedztwa oparte na haszowaniu przestrzennym
          z~reorganizacją buforów, zapewniające stały czas zapytania
          i~skoalescowane odczyty pamięci GPU;
    \item zaimplementowano i~porównano ilościowo dwa algorytmy sortowania
          na GPU --- sortowanie bitoniczne i~pozycyjne --- wykazując niemal
          sześciokrotną przewagę tego drugiego w~pełnym potoku symulacji;
    \item zaprojektowano organizację danych SOA z~podwójnym buforowaniem pozycji
          i~prędkości, strukturalnie eliminującą wyścigi danych między wątkami
          GPU;
    \item zbudowano wolumetryczny potok wizualizacji: splatting gęstości
          do tekstury trójwymiarowej, przemarsz promienia z~bisekcją oraz model
          optyczny wody łączący efekt Fresnela, absorpcję Beer-Lamberta,
          rozproszenie podpowierzchniowe, odbłysk GGX i~oświetlenie panoramą
          HDRI;
    \item przeprowadzono analizę zbieżności obu solverów oraz eksperymenty
          z~adaptacyjnym krokiem czasowym według warunku CFL, dokumentując
          źródła ich ograniczeń;
    \item udostępniono interaktywny panel konfiguracji umożliwiający modyfikację
          parametrów fizycznych i~wizualnych w~trakcie działania symulacji.
\end{itemize}

Wynikiem pracy jest działający symulator zdolny do stabilnej animacji dziesiątek
tysięcy cząstek w~czasie rzeczywistym, wraz z~zestawem narzędzi diagnostycznych
i~pomiarowych. Przeprowadzone testy ujawniły zarazem najważniejsze ograniczenie
przyjętego modelu: brak właściwej obsługi brzegu powoduje trwały deficyt
gęstości przy swobodnej powierzchni, który blokuje kryterium zbieżności solvera
i~uniemożliwia bezpieczne stosowanie adaptacyjnego kroku czasowego
(podrozdz.~\ref{subsec:cfl}). Świadomy powrót do stałego $\Delta t$ i~stałej
liczby iteracji stanowi udokumentowany kompromis między stabilnością
a~wydajnością.

Niniejsza praca koncentrowała się przede wszystkim na warstwie fizycznej
symulacji; ta orientacja wyznacza naturalne kierunki dalszego rozwoju systemu.

\begin{itemize}
    \item \textbf{Uzupełnienie modelu fizycznego.} Pierwszym krokiem powinna być
          obsługa brzegu za pomocą cząstek-duchów lub stałych warstw cząstek
          brzegowych~\cite{schechter2012ghost}, która usunęłaby deficyt jądra
          przy powierzchni i~odblokowała zarówno kryterium
          $\varepsilon$-zbieżności, jak i~adaptacyjny warunek CFL. Kolejnym
          naturalnym rozszerzeniem jest implementacja napięcia powierzchniowego,
          otwierająca drogę do zjawisk kapilarnych, formowania kropel
          i~realistycznego zachowania cienkich warstw cieczy.
    \item \textbf{Rozbudowa warstwy wizualizacji.} Ponieważ wizualizacja pełniła
        w~pracy rolę służebną wobec fizyki, pozostaje tu szerokie pole
        eksperymentów. Obejmują one jawną rekonstrukcję powierzchni metodą
        maszerujących sześcianów (ang.~\textit{marching cubes}), techniki
        renderowania w~przestrzeni ekranu oraz cząstki wtórne piany i~rozbryzgów.
        Dodatkowo, wdrożenie sprzętowego śledzenia promieni poprzez rozszerzenie
        \texttt{VK\_KHR\_ray\_tracing\_pipeline} pozwoliłoby na uzyskanie fizycznie
        poprawnych kaustyk i~wielokrotnych odbić światła w~objętości cieczy.
    \item \textbf{Techniki optymalizacji.} Obiecującym kierunkiem są
          skompresowane listy sąsiedztwa~\cite{band2019compressed}, których
          pełna realizacja na GPU pozostaje otwartym problemem badawczym,
          a~ponadto adaptacyjna rozdzielczość cząstek zagęszczająca symulację
          wyłącznie w~obszarach istotnych wizualnie oraz głębsze wykorzystanie
          pamięci dzielonej i~operacji podgrupowych
          (ang.~\textit{subgroup operations}) w~shaderach solvera.
    \item \textbf{Integracja z~ekosystemem.} Silnik może zostać przekształcony
          w~bibliotekę wielokrotnego użytku lub rozwinięty w~stronę
          pełnoprawnego silnika aplikacji interaktywnych; alternatywą jest
          opakowanie go we wtyczkę do istniejącego środowiska, analogicznie
          do roli, jaką NiagaraFluids pełni w~Unreal Engine~5. Krokiem
          poprzedzającym byłoby usunięcie zależności od platformy Windows
          i~udostępnienie wersji dla systemów Linux i~macOS.
    \item \textbf{Metody oparte na danych.} Odrębną perspektywę otwiera
          częściowe lub całkowite zastąpienie obliczeń fizycznych modelami
          statystycznymi i~sieciami neuronowymi: aproksymacja kroku solvera
          modelem nauczonym na danych z~symulacji referencyjnych, generowanie
          detali wysokiej częstotliwości na podstawie zgrubnej symulacji
          o~małej liczbie cząstek czy statystyczna analiza pól prędkości
          i~gęstości pozwalająca przewidywać zachowanie cieczy bez jawnego
          rozwiązywania równań ruchu.
\end{itemize}

Przeprowadzone prace potwierdzają, że połączenie języka Rust i~API Vulkan
stanowi dojrzały fundament dla wysokowydajnych symulacji fizycznych: statyczne
gwarancje kompilatora wyeliminowały całą klasę błędów zarządzania pamięcią,
a~jawna kontrola nad procesorem graficznym pozwoliła przenieść na niego
kompletny potok obliczeniowy. Powstały system --- wraz z~opisanymi w~pracy
pomiarami, diagnozami i~zidentyfikowanymi ograniczeniami --- może służyć jako
platforma odniesienia i~punkt wyjścia dla wymienionych wyżej kierunków badań.
